% GENERATED FILE. Edit canonical lessons, then run npm run generate:books.
% canonical-source-hash: fnv1a64:f2120f21c95a5171
% canonical-lessons: LA-C316-eripio, LA-C316-incolumis, LA-C316-subvenio, LA-C316-succurro, LA-C316-auxilium

\chapter{Auxilium}
\label{ch:la-auxilium}
\hlchaptermodality{316}

\begin{chapteropening}
\textbf{By the end of this chapter:} \emph{I can tell how people were rescued: to snatch away, unharmed, to come to help, to run to help and help.}

Complete the last lesson of chapter 316: I can tell how people were rescued: to snatch away, unharmed, to come to help, to run to help and help.
\end{chapteropening}

\section[ēripiō, ēripere]{ēripiō, ēripere — to snatch away, to rescue}

\label{lesson:LA-C316-eripio}



\begin{quote}
Before the new one: say the Latin for to overwhelm, then the Latin for a disaster.
\end{quote}

\subsection*{You'll want to know: ēripiō, ēripere}
\textbf{ēripiō, ēripere} — \textquotedblleft{}to snatch away, to rescue\textquotedblright{}.

\begin{quote}
Pater fīlium ē flūmine ēripuit.
\end{quote}

\begin{quote}
\emph{The father snatched his son out of the river.}
\end{quote}

\emph{ēripiō, ēripere, ēripuī, ēreptum}: \emph{ē} + \emph{rapiō}. \emph{Ēripere ex perīculō}: to rescue from danger.

\begin{cousinweb}
\textbf{ēripiō} is built on \emph{rapiō}, \textquotedblleft{}to snatch\textquotedblright{}, with \emph{ex-}, \textquotedblleft{}out\textquotedblright{} (EtymDB); the vowel changes inside the compound.
\end{cousinweb}

\begin{grammarlens}[title={What you've built}]
One more word for telling, explaining and describing at B1.
\end{grammarlens}

\subsection*{Your turn}
\begin{itemize}
\raggedright
  \item \emph{Say it:} \emph{ēripiō, ēripere}
  \item \emph{Say it:} \emph{ēripiō, ēripere}, once more
  \item \emph{Say it:} a sentence of your own with \emph{ēripiō, ēripere}
  \item \emph{Recall:} say the Latin for to overwhelm, then the Latin for a disaster, then say \emph{ēripiō, ēripere} again
\end{itemize}

\begin{tcolorbox}[breakable,colback=teal!4,colframe=teal!35!black,title={Before you move on}]
What does ēripiō, ēripere mean? (\textquotedblleft{}To snatch away, to rescue\textquotedblright{}.) Say it once more.
\end{tcolorbox}
\section[incolumis, incolume]{incolumis, incolume — unharmed, safe}

\label{lesson:LA-C316-incolumis}



\begin{quote}
Before the new one: say the Latin for a disaster, then the Latin for to snatch away.
\end{quote}

\subsection*{You'll want to know: incolumis, incolume}
\textbf{incolumis, incolume} — \textquotedblleft{}unharmed, safe\textquotedblright{}.

\begin{quote}
Omnēs nautae incolumēs ad lītus pervēnērunt.
\end{quote}

\begin{quote}
\emph{All the sailors reached the shore unharmed.}
\end{quote}

Two endings. Caesar uses it of soldiers who came back safe and sound.

\begin{cousinweb}
\textbf{incolumis} joins \emph{in-}, \textquotedblleft{}not\textquotedblright{}, and a second part that EtymDB traces to Proto-Italic \emph{*kadami-}, \textquotedblleft{}damaged\textquotedblright{}. Lewis \& Short compare \emph{calamitās}; the details are not settled.
\end{cousinweb}

\begin{grammarlens}[title={What you've built}]
One more word for telling, explaining and describing at B1.
\end{grammarlens}

\subsection*{Your turn}
\begin{itemize}
\raggedright
  \item \emph{Say it:} \emph{incolumis, incolume}
  \item \emph{Say it:} \emph{incolumis, incolume}, once more
  \item \emph{Say it:} a sentence of your own with \emph{incolumis, incolume}
  \item \emph{Recall:} say the Latin for a disaster, then the Latin for to snatch away, then say \emph{incolumis, incolume} again
\end{itemize}

\begin{tcolorbox}[breakable,colback=teal!4,colframe=teal!35!black,title={Before you move on}]
What does incolumis, incolume mean? (\textquotedblleft{}Unharmed, safe\textquotedblright{}.) Say it once more.
\end{tcolorbox}
\section[subveniō, subvenīre]{subveniō, subvenīre — to come to help}

\label{lesson:LA-C316-subvenio}



\begin{quote}
Before the new one: say the Latin for to snatch away, then the Latin for unharmed.
\end{quote}

\subsection*{You'll want to know: subveniō, subvenīre}
\textbf{subveniō, subvenīre} — \textquotedblleft{}to come to help\textquotedblright{}.

\begin{quote}
Vīcīnī celeriter subvēnērunt.
\end{quote}

\begin{quote}
\emph{The neighbours quickly came to help.}
\end{quote}

\emph{subveniō, subvenīre, subvēnī, subventum}, with the dative: \emph{mihi subvēnit}. English \emph{subvention}.

\begin{cousinweb}
\textbf{subveniō} is built on \emph{veniō}, \textquotedblleft{}to come\textquotedblright{}, with \emph{sub-} (EtymDB): to come up in support, as reserves do in battle. English \emph{subvention} comes from it.
\end{cousinweb}

\begin{grammarlens}[title={What you've built}]
One more word for telling, explaining and describing at B1.
\end{grammarlens}

\subsection*{Your turn}
\begin{itemize}
\raggedright
  \item \emph{Say it:} \emph{subveniō, subvenīre}
  \item \emph{Say it:} \emph{subveniō, subvenīre}, once more
  \item \emph{Say it:} a sentence of your own with \emph{subveniō, subvenīre}
  \item \emph{Recall:} say the Latin for to snatch away, then the Latin for unharmed, then say \emph{subveniō, subvenīre} again
\end{itemize}

\begin{tcolorbox}[breakable,colback=teal!4,colframe=teal!35!black,title={Before you move on}]
What does subveniō, subvenīre mean? (\textquotedblleft{}To come to help\textquotedblright{}.) Say it once more.
\end{tcolorbox}
\section[succurrō, succurrere]{succurrō, succurrere — to run to help}

\label{lesson:LA-C316-succurro}



\begin{quote}
Before the new one: say the Latin for unharmed, then the Latin for to come to help.
\end{quote}

\subsection*{You'll want to know: succurrō, succurrere}
\textbf{succurrō, succurrere} — \textquotedblleft{}to run to help\textquotedblright{}.

\begin{quote}
Puellae in aquā clāmantī duo nautae succurrērunt.
\end{quote}

\begin{quote}
\emph{Two sailors ran to help the girl who was shouting in the water.}
\end{quote}

\emph{succurrō, succurrere, succurrī, succursum}, with the dative. English \emph{succour}.

\begin{cousinweb}
\textbf{succurrō} is built on \emph{currō}, \textquotedblleft{}to run\textquotedblright{}, with \emph{sub-} (EtymDB): to run up to someone. English \emph{succour} comes from it through French.
\end{cousinweb}

\begin{grammarlens}[title={What you've built}]
One more word for telling, explaining and describing at B1.
\end{grammarlens}

\subsection*{Your turn}
\begin{itemize}
\raggedright
  \item \emph{Say it:} \emph{succurrō, succurrere}
  \item \emph{Say it:} \emph{succurrō, succurrere}, once more
  \item \emph{Say it:} a sentence of your own with \emph{succurrō, succurrere}
  \item \emph{Recall:} say the Latin for unharmed, then the Latin for to come to help, then say \emph{succurrō, succurrere} again
\end{itemize}

\begin{tcolorbox}[breakable,colback=teal!4,colframe=teal!35!black,title={Before you move on}]
What does succurrō, succurrere mean? (\textquotedblleft{}To run to help\textquotedblright{}.) Say it once more.
\end{tcolorbox}
\section[auxilium, auxiliī]{auxilium, auxiliī — help, aid}

\label{lesson:LA-C316-auxilium}



\begin{quote}
Before the new one: say the Latin for to come to help, then the Latin for to run to help.
\end{quote}

\subsection*{You'll want to know: auxilium, auxiliī}
\textbf{auxilium, auxiliī} — \textquotedblleft{}help, aid\textquotedblright{}.

\begin{quote}
Fēmina auxilium ā vīcīnīs petīvit.
\end{quote}

\begin{quote}
\emph{The woman asked the neighbours for help.}
\end{quote}

Neuter. The plural \emph{auxilia} means auxiliary troops. English \emph{auxiliary}.

\begin{cousinweb}
\textbf{auxilium} is built on \emph{augeō}, \textquotedblleft{}to increase\textquotedblright{} (EtymDB): what adds to your strength. English \emph{auxiliary} comes from it.
\end{cousinweb}

\begin{grammarlens}[title={What you've built}]
One more word for telling, explaining and describing at B1.
\end{grammarlens}

\subsection*{Your turn}
\begin{itemize}
\raggedright
  \item \emph{Say it:} \emph{auxilium, auxiliī}
  \item \emph{Say it:} \emph{auxilium, auxiliī}, once more
  \item \emph{Say it:} a sentence of your own with \emph{auxilium, auxiliī}
  \item \emph{Recall:} say the Latin for to come to help, then the Latin for to run to help, then say \emph{auxilium, auxiliī} again
\end{itemize}

\begin{tcolorbox}[breakable,colback=teal!4,colframe=teal!35!black,title={Before you move on}]
What does auxilium, auxiliī mean? (\textquotedblleft{}Help, aid\textquotedblright{}.) Say it once more.
\end{tcolorbox}
